\section{Introduction}
\label{sec:intro}

Learned surrogates of dynamical systems are usually trained to predict one step ahead and then used autoregressively, with each prediction fed back as the next input. In this mode of operation, the approximation error made at one step becomes a perturbation of all subsequent states, and a surrogate with a small one-step error can still produce large errors over long rollouts \citep{bengio2015,venkatraman2015,lambert2022}. For this reason, neural weather models and neural PDE solvers are evaluated over many autoregressive steps and are often fine-tuned with multistep objectives \citep{brandstetter2022,lam2023,list2025}.

The growth of a perturbation along a trajectory is governed by the tangent dynamics, and several works use Jacobian information of the learned map to explain or control long-horizon errors, either by regularizing one-step Jacobians \citep{pan2018,nie2026} or by penalizing the non-normality and non-commutativity of consecutive Jacobians, which control the transient amplification available to perturbations \citep{pervez2026}. The leading singular value of a finite-time propagator, however, describes the worst perturbation. The rollout error of a surrogate initialized on the correct state is generated by the defects that the surrogate makes at every step, and whether these defects excite the most amplified directions, and whether their transported contributions reinforce or cancel each other, is not determined by the propagator alone.

The accounting of how local errors are transported and summed is classical in numerical analysis, where the global error of a time integrator is the sum of transported local errors \citep[the ``Lady Windermere's fan'' of][]{hairer1993}, and in predictability studies of model error \citep{nicolis2003,vannitsem2002}. A learned surrogate can be regarded as an inexpensive and imperfect numerical scheme whose local error is the approximation error of the network. Unlike the local error of a classical scheme, this error is not controlled by a truncation order. It is a deterministic and largely unknown function of the state, shaped by the training data and by the architecture, and its interaction with the finite-time geometry of the learned dynamics has to be measured.

The current work quantifies how much of the amplification available in the learned dynamics is realized by the defects of the surrogate. The first-order rollout error is written as a sum of defects transported by reference-trajectory propagators, and an exact factorization splits the rollout error into five factors: the magnitude of the injected defects, the worst-case gain of the tail propagators, the efficiency with which the defects use that gain, the cancellation among the transported defects and the nonlinear remainder. Since the factorization is exact in logarithmic form, the variance of the log-error across initial states is attributed to the five factors without approximation, and the same attribution applies to the difference between two surrogates evaluated on the same forecasts. The conclusion that the realized error depends on the actual defect field, and not on its statistics, is then used to build an error estimate that needs no reference data.

The analysis is applied to residual multilayer perceptrons (MLPs) trained on the 20-dimensional Lorenz--96 system with three optimization budgets, and to Fourier neural operators (FNOs) trained on the Kuramoto--Sivashinsky equation. Horizons are expressed in Lyapunov times, so that the two systems are compared over the same range of chaotic error growth. The main contributions of the current work are the following.
\begin{enumerate}
\item An exact five-factor factorization of the autoregressive rollout error, null models for the efficiency (isotropic and covariance-matched) and for the vector composition (coherent, incoherent and bias-only limits), and an exact variance attribution of the log-error (Section~\ref{sec:accounting}).
\item Evidence, for two systems, two architectures and a sixtyfold range of one-step accuracy, that the defects realize, within 7\%, the fraction of the worst-case gain expected for random vectors with their own covariance. The decrease of the efficiency with the horizon follows from the singular spectrum of the propagators (Section~\ref{sec:res_efficiency}).
\item Evidence that the transported defects are strongly coherent, about four times more than independent defects, and that the coherence comes from the temporal correlation of the defects along the trajectory and not from a systematic bias (Section~\ref{sec:res_coherence}).
\item A variance attribution showing that the defect magnitude explains the differences between forecasts at short horizons, and that the worst-case gain, the coherence and, for the FNOs, the orientation of the defects take over beyond about half a Lyapunov time, with a similar budget for all accuracy levels (Section~\ref{sec:res_budget}).
\item A reference-free error estimate in which the disagreement between independently trained surrogates replaces the unknown defects. Except for the under-trained Lorenz--96 surrogates, it ranks individual forecasts up to about one Lyapunov time and is calibrated in magnitude for the converged ones, whereas scores built from the defect statistics do not rank the forecasts (Section~\ref{sec:res_ensemble}).
\item A diagnosis of multistep fine-tuning. It reduces the error of under-trained surrogates, whose free rollouts leave the attractor, by lowering the gain of their tangent dynamics, and it has no consistent effect on surrogates that reproduce the reference tangent dynamics (Sections~\ref{sec:res_fidelity} and \ref{sec:res_intervention}).
\end{enumerate}
